\chapter{Introdução}
\label{chap:intro} %este label será usado para referenciar este capítulo

% As primeiras frases têm a missão de prender a atenção do leitor e por isso são as mais importantes do texto. Diga o quanto antes o que você fez e quais são os resultados alcançados. Ao terminar de ler a introdução o leitor tomará uma nova decisão de se vale a pena ou não continuar lendo o texto. Capte a atenção do leitor bem aqui.

% A comunicação escrita é considerada umas das cinco habilidades mais importantes por profissionais de engenharia e um engenheiro passa em média mais de $25$\% do seu tempo escrevendo \cite{eggert2002response,spretnak1982survey}. Uma quantidade similar de tempo é gasta na escrita de correspondência e de relatórios técnicos \cite{cunningham2012perceptions}. Dessa forma, encare a escrita do seu projeto como um treinamento nessa importante habilidade.

% Neste texto você encontrará não apenas uma estrutura para escrever seu trabalho em \LaTeX, mas também um pequeno manual de boas práticas na escrita técnica. Leia com atenção e coloque as sugestões em prática à medida que preenche o texto com o conteúdo do seu próprio projeto. Também será apresentado um número de vícios de escrita comumente encontrados nas monografias de alunos. 

% A seguir está a estrutura de organização sugerida pelo colegiado do curso. Note que ela não é necessariamente a melhor para contar a história do seu projeto. Você pode por exemplo preferir usar títulos mais pertinentes ao seu contexto. Contudo, o seu texto deve conter cada um dos pontos a seguir.

\section{Motivação e Justificativa}
\label{sec:motivacao}

A gestão eficiente do consumo de energia elétrica em ambientes residenciais representa um desafio contemporâneo crucial, impulsionado tanto pela crescente demanda energética global quanto pela urgência na adoção de práticas sustentáveis. Frequentemente, os usuários domésticos carecem de visibilidade detalhada sobre seus perfis de consumo e de ferramentas acessíveis para o gerenciamento ativo, o que cria obstáculos para a identificação de desperdícios e a adoção de medidas eficazes para a otimização energética e redução de custos. Essa lacuna é exacerbada pela escassez de sistemas de monitoramento e controle que sejam simultaneamente de fácil implementação em residências convencionais e intuitivos para o usuário final, dificultando a transição para um uso mais consciente e econômico da energia no cotidiano.

Paralelamente, o campo da automação residencial, intrinsecamente ligado à gestão energética inteligente, confronta-se com obstáculos significativos de interoperabilidade. A heterogeneidade de dispositivos, sensores e tecnologias de comunicação, provenientes de múltiplos fabricantes e frequentemente baseados em padrões proprietários ou concorrentes, resulta em ecossistemas fragmentados. Essa fragmentação impede a comunicação eficiente entre componentes e a consequente criação de sistemas de monitoramento e controle verdadeiramente unificados e coesos. Tal cenário não apenas subutiliza o potencial da automação para otimizar o consumo energético e elevar a qualidade de vida, mas também restringe o acesso a soluções tecnológicas que poderiam fomentar ambientes domésticos mais confortáveis, seguros e ambientalmente responsáveis.

Agrava este panorama o fato de que uma vasta parcela do parque imobiliário residencial existente não foi concebida com infraestrutura para automação. A modernização desses lares para incorporar tecnologias de monitoramento e controle energético inteligente frequentemente implica intervenções estruturais complexas e custos proibitivos para muitos proprietários. Essa barreira à adoção, tanto técnica quanto financeira, priva um número expressivo de residências dos benefícios tangíveis da automação — como maior eficiência energética, conforto aprimorado e segurança reforçada. Consequentemente, a carência de soluções que sejam simultaneamente acessíveis, de fácil instalação em edificações convencionais e eficazes na gestão do consumo energético configura uma oportunidade premente para o desenvolvimento de inovações capazes de transpor esses obstáculos e democratizar o acesso a um gerenciamento energético residencial mais inteligente e sustentável.

\section{Objetivos do Projeto}
\label{sec:objetivos}

Tendo em vista o exposto acima, este projeto tem por objetivos:

\begin{enumerate}[a)]
\item Realizar uma pesquisa e análise de tecnologias existentes que permitam o monitoramento energético sem a necessidade de infraestrutura prévia.
\item Explorar soluções adequadas para residências convencionais, que não foram originalmente projetadas para automação residencial.
\item Desenvolver um protótipo de sistema de monitoramento e controle de fácil instalação e uso, focado na aplicação de princípios de controle estatístico de processo.
\item Implementar algoritmos para processamento de dados, previsão de consumo e aplicação de técnicas de controle estatístico de processo.
\item Desenvolver uma plataforma de software intuitiva para visualização, análise dos dados e suporte ao controle estatístico de processo pelo usuário.
\item Testar e validar o sistema em ambiente real, garantindo sua eficiência e usabilidade.
\item Promover a viabilidade de implementação em larga escala por meio da análise de custos e estratégias de adoção.
\item Documentar e divulgar os resultados obtidos, destacando as contribuições para a Engenharia de Controle e Automação e os benefícios sociais e ambientais proporcionados pelo sistema desenvolvido.
\end{enumerate}

\section{Local de Realização}
\label{sec:empresa}

O local de realização deste projeto é a residência do próprio desenvolvedor, onde os testes dos equipamentos e a simulação das necessidades reais de automação residencial estão sendo realizados. A escolha deste ambiente se justifica pela possibilidade de realizar uma avaliação prática e realista das soluções de monitoramento e controle estatístico de processos para o consumo energético, simulando as condições do dia a dia em uma residência convencional. A residência não foi originalmente projetada com infraestrutura para automação, o que representa um desafio adicional, mas também um aspecto relevante para a validação do sistema, pois permite testar a adaptação de tecnologias em um ambiente que não conta com recursos avançados de automação.

Neste contexto, o projeto visa avaliar a viabilidade de implementação de soluções de monitoramento energético em ambientes residenciais que, assim como muitas casas, não foram concebidos para suportar sistemas automatizados. Dessa forma, o projeto busca não apenas desenvolver uma solução eficiente, mas também uma que seja de fácil adaptação e acessível para a grande maioria dos consumidores. O vínculo com o local de testes é direto, pois, além de ser o ambiente de realização das simulações, a residência do desenvolvedor também serve como um campo de experimentação e validação para os resultados do sistema proposto.

\section{Estrutura da Monografia}
\label{sec:organizacao}

O trabalho está dividido em cinco capítulos. Neste primeiro capítulo (Capítulo \ref{chap:intro}) apresentou uma introdução ao projeto, incluindo a motivação, os objetivos e o local de realização. No Capítulo \ref{chap:principios} (Fundamentos da Automação Residencial e Monitoramento Energético) descrevem-se os princípios básicos e modelos arquiteturais para sistemas de automação residencial e monitoramento energético, abordando as camadas de aquisição de dados, processamento, gerenciamento, apresentação e automação. No Capítulo 3 (Metodologia) explora-se a metodologia de desenvolvimento do protótipo proposto, incluindo a pesquisa e análise de tecnologias, a seleção de ferramentas de software, o desenvolvimento do protótipo com configuração de sensores, a implementação da pipeline de dados com Dagster, o desenvolvimento da interface de usuário com Streamlit, e a implementação de algoritmos de processamento e análise, como o Controle Estatístico de Processos (CEP). No Capítulo 4 (Resultados) apresentam-se os resultados da implementação do sistema, detalhando o pipeline de dados e a aplicação de interface do usuário com suas funcionalidades e visualizações. Finalmente, no Capítulo \ref{sec:conclusoes} (Conclusões) resumem-se os resultados alcançados, discute-se o atendimento aos objetivos, as principais contribuições, limitações e propõem-se trabalhos futuros.

\clearpage
